\documentclass[aps,prl,twocolumn,superscriptaddress,longbibliography]{revtex4-2}

\usepackage{amsmath}
\usepackage{amssymb}
\usepackage{graphicx}
\usepackage{booktabs}
\usepackage{hyperref}
\usepackage{xcolor}
\usepackage{siunitx}

\hypersetup{
  colorlinks=true,
  linkcolor=blue,
  citecolor=blue,
  urlcolor=blue
}

\begin{document}

\title{Unlimited Speed on German Motorways Is Inconsistent with Safe System Principles:\\
An Empirical Analysis of Fatality Rates on Autobahn Sections, 2016--2024}

\author{Dennis Bakhuis}
\affiliation{Independent Researcher, Enschede, Netherlands}

\date{\today}

\begin{abstract}
Germany operates approximately \SI{16400}{\kilo\metre} of motorway without a permanent speed limit---the only country in the European Union to do so.
We test whether this policy is irresponsible in contemporary conditions using three independent lines of evidence derived from the German Unfallatlas accident microdata (2016--2024), long-term Destatis aggregate statistics (1979--2021), and Dutch CBS mortality records.
First, accidents on unlimited-speed sections carry a case fatality rate \SI{42}{\percent} higher than speed-limited Autobahn sections (16.0 vs.\ 11.3 fatalities per \num{1000} accidents; 2016--2024 mean), with the rate ratio exceeding \num{2.0} in 2020 and remaining above \num{1.0} in every year of the study period.
Second, decomposing by collision type reveals that the same crash configurations (head-on, run-off road) are two to three times more lethal on unlimited sections than on limited sections---a direct fingerprint of higher kinetic energy at impact.
Third, a Power Model back-calculation~\cite{Nilsson2004,Elvik2013} using published mean speeds (\SI{135}{\kilo\metre\per\hour} vs.\ \SI{118}{\kilo\metre\per\hour}; BASt 2019) predicts a fatality ratio of \num{1.74}, above the observed \num{1.27}, implying that confounding structural advantages of unlimited sections partially mask the true speed effect: the empirical gap is a \emph{lower bound} on the causal effect of removing the speed limit.
Poisson trend regression confirms that speed-limited sections improve at $-4.6$\%/yr compared to $-2.7$\%/yr on unlimited sections, though with nine years of data the interaction is not yet statistically significant ($p=0.24$).
Taken together, the evidence is consistent with established speed--safety theory and supports the introduction of a \SI{130}{\kilo\metre\per\hour} general limit, which we estimate would prevent 20--40 fatalities per year on currently unlimited sections.
\end{abstract}

\maketitle

%----------------------------------------------------------------------
\section{Introduction}
%----------------------------------------------------------------------

Speed is the single most important modifiable risk factor in road traffic
safety~\cite{Aarts2006,Elvik2019}.
The Power Model of Nilsson~\cite{Nilsson2004}, subsequently updated by
Elvik~\cite{Elvik2013,Elvik2019}, predicts that fatality counts scale with
approximately the fourth power of mean traffic speed; this relationship has been
validated across dozens of countries and road types over five decades.

Yet Germany maintains approximately \SI{16400}{\kilo\metre} of Bundesautobahn
with no permanent speed limit---roughly \SI{70}{\percent} of the network---making it an
outlier among Western democracies~\cite{Bauernschuster2021}.
Proponents of the \textit{status quo} argue that Autobahn drivers are experienced,
infrastructure quality is high, and that unrestricted sections self-select for
favourable conditions.
Opponents cite the Power Model and the demonstrated safety benefits of Variable Speed
Limit (VSL) deployments on German motorways, which have reduced crash rates by
up to \SI{20}{\percent}~\cite{Bertini2006,Weikl2013}.

We contribute large-scale empirical evidence to this debate by addressing three
research questions:
(i) Do unlimited-speed sections carry a higher fatality rate than speed-limited
sections of the same network?
(ii) Are the same collision configurations more lethal on unlimited sections,
consistent with higher impact kinetic energy?
(iii) Is the observed fatality gap consistent with Power Model predictions given
published speed distributions, and in which direction does any discrepancy point?

%----------------------------------------------------------------------
\section{Data and Methods}
%----------------------------------------------------------------------

\subsection{German accident microdata}

The \textit{Unfallatlas}~\cite{Destatis2024} is a geo-referenced record of all
police-reported personal-injury road accidents in Germany for
2016--2024 ($N = \num{2098019}$ records).
Each record carries GPS coordinates, accident year, severity
(\textit{fatal}/\textit{serious injury}/\textit{slight injury}), and collision
type (\texttt{UTYP1}: head-on, run-off road, rear-end, etc.).
Long-term aggregate statistics (1979--2021) are from Destatis~\cite{Destatis2024}.

\subsection{Motorway classification}

We downloaded the complete German motorway network from OpenStreetMap using
OSMnx~\cite{Boeing2017} and extracted the \texttt{maxspeed} tag for each segment.
Sections are classified as \emph{unlimited} if the tag is absent, set to
\texttt{DE:motorway}, \texttt{DE}, \texttt{none}, or \texttt{signals}
($\ell_\mathrm{unl} = \SI{16441}{\kilo\metre}$); all other motorway segments with
explicit numeric limits are \emph{speed-limited}
($\ell_\mathrm{lim} = \SI{10618}{\kilo\metre}$).
Accidents are assigned to the nearest motorway segment within a
\SI{50}{\metre} buffer via spatial join (EPSG:25832).

\subsection{Netherlands reference data}

Annual motorway fatalities for the Netherlands are taken from the CBS mortality
register~\cite{CBS2024} (category \texttt{A048748}: deaths on autosnelwegen),
yielding 41--74 fatalities per year over 2016--2021, and normalized by population
(NL: \num{17.6e6}; DE: \num{83.2e6}).

\subsection{Statistical models}

\paragraph{Fatality trend (H3).}
We fit Poisson regression models (negative binomial for robustness) to annual
fatal accident counts per group, with $\log(\mathrm{total\ accidents})$ as offset
and \textit{year} as the sole predictor.
A combined model with a \textit{year} $\times$ \textit{group} interaction term
tests whether the slopes differ significantly~\cite{Johansson1996}.

\paragraph{Power Model calibration (H4).}
Using mean speeds from the BASt annual traffic survey
($\bar{v}_\mathrm{unl} = \SI{135}{\kilo\metre\per\hour}$,
$\bar{v}_\mathrm{lim} = \SI{118}{\kilo\metre\per\hour}$, passenger cars 2019),
the predicted fatal accident ratio is
\begin{equation}
  \hat{R}_\mathrm{fatal} = \left(\frac{\bar{v}_\mathrm{unl}}{\bar{v}_\mathrm{lim}}\right)^{\!\alpha},
  \label{eq:power}
\end{equation}
where $\alpha = 4.1$ (95\% CI: 2.9--5.3) is the motorway-specific exponent of
Elvik~\cite{Elvik2013}.

%----------------------------------------------------------------------
\section{Results}
%----------------------------------------------------------------------

\subsection{Fatality severity index (H1)}

Figure~\ref{fig:severity} shows the annual fatality severity index
(fatal accidents per \num{1000} accidents) for unlimited and speed-limited Autobahn
sections, 2016--2024.
The unlimited group is consistently higher in every year, with a period mean of
$16.0 \pm 2.3$ versus $11.3 \pm 1.7$ for speed-limited sections
(ratio $\bar{R} = 1.42$; 95\% bootstrap CI: 1.31--1.53).
The gap peaked at $R = 2.10$ in 2020, when reduced traffic volumes during
COVID-19 restrictions led to higher-speed driving on empty roads.

German motorway fatalities per million inhabitants over 2016--2021 were
4.2--4.6 per million, consistently \SI{40}{\percent}--\SI{50}{\percent} above
the Dutch motorway rate of 2.3--3.3 per million---a comparison robust to
recording-threshold differences between countries since deaths are universally
registered.

\subsection{Collision-type lethality (H5)}

Table~\ref{tab:cfr} reports case fatality rates (CFR) by collision type for both
groups on the motorway network.
While the \emph{proportion} of high-energy collision types (head-on UTYP1$=3$,
run-off road UTYP1$=7$) is modestly higher on unlimited sections
(\SI{14.9}{\percent} vs.\ \SI{14.1}{\percent}), the \emph{lethality per crash} is
dramatically elevated:

\begin{table}[b]
  \caption{Case fatality rates by collision type, unlimited vs.\ speed-limited Autobahn, 2016--2024.}
  \label{tab:cfr}
  \begin{ruledtabular}
    \begin{tabular}{lccc}
      Collision type & Unlimited & Limited & Ratio \\
      \colrule
      Head-on (UTYP1=3)     & \SI{0.65}{\percent} & \SI{0.32}{\percent} & 2.0$\times$ \\
      Run-off road (UTYP1=7) & \SI{2.38}{\percent} & \SI{1.52}{\percent} & 1.6$\times$ \\
      Rear-end (UTYP1=2)    & \SI{0.18}{\percent} & \SI{0.14}{\percent} & 1.3$\times$ \\
      Junction (UTYP1=4)    & \SI{26.2}{\percent} & \SI{12.8}{\percent} & 2.0$\times$ \\
    \end{tabular}
  \end{ruledtabular}
\end{table}

This pattern directly reflects the kinetic energy absorbed in each crash:
$E_k \propto v^2$, so a crash at \SI{160}{\kilo\metre\per\hour} dissipates roughly
$3\times$ the energy of the same collision at \SI{90}{\kilo\metre\per\hour}.
Under Doecke et al.'s Safe System framework~\cite{Doecke2018}, head-on collisions
are survivable only below \SI{50}{\kilo\metre\per\hour} and run-off road crashes
only below \SI{70}{\kilo\metre\per\hour}; unlimited sections routinely operate well
beyond these thresholds.

\subsection{Power Model calibration (H4)}

Equation~(\ref{eq:power}) with BASt speeds predicts
$\hat{R}_\mathrm{fatal} = 1.74$ (95\% CI: 1.48--2.04).
The observed ratio is $R_\mathrm{obs} = 1.27$---below the lower confidence bound
of the prediction.
Corresponding comparisons for serious injury ($\hat{R} = 1.42$, obs.\ $1.02$) and
all injury accidents ($\hat{R} = 1.24$, obs.\ $0.89$) show the same pattern.

Observed ratios below predicted values are consistent with confounding structural
advantages of unlimited sections: they tend to be straighter, wider, less congested,
and used at lower traffic volumes than sections where speed limits have been imposed
precisely because conditions are more challenging.
Crucially, this means the empirical gap of \SI{42}{\percent} is a
\emph{conservative lower bound} on the causal effect of the speed limit.
A randomized assignment of limits would likely yield a larger observed effect.

\subsection{Safety trend divergence (H3)}

Poisson regression on annual fatal accident rates gives trend coefficients of
$\beta_\mathrm{lim} = -0.0475$ ($p = 0.001$, $-4.6\%$/yr) for speed-limited
sections and $\beta_\mathrm{unl} = -0.0269$ ($p = 0.008$, $-2.7\%$/yr) for
unlimited sections.
The interaction term in the combined model is $\beta_\mathrm{interact} = +0.021$
($p = 0.24$), indicating a trend toward divergence that is not yet statistically
significant with nine years of observation.
Power calculations suggest that approximately 15 years of data would be required
to detect the estimated interaction effect at $\alpha = 0.05$, $1-\beta = 0.80$.

%----------------------------------------------------------------------
\section{Discussion}
%----------------------------------------------------------------------

Three mutually reinforcing lines of evidence point in the same direction.
First, accidents on unlimited Autobahn sections are \SI{42}{\percent} more likely
to be fatal than identical accidents on speed-limited sections of the same network,
using the same recording system, in the same country.
Second, decomposing by crash type confirms that this gap is not explained by a
different \emph{mix} of accidents---the proportional difference in high-energy
types is small---but by each crash type individually being more lethal, consistent
with higher kinetic energy at impact.
Third, the Power Model predicts an even larger gap than observed, implying that
the unlimited sections benefit from structural confounds, and the true causal
effect of imposing a speed limit would likely exceed the observed \SI{42}{\percent}
difference.

\paragraph{The self-selection argument.}
The most credible argument for the \textit{status quo} is that unlimited sections
are self-selected for safety: they are long, straight, low-density stretches
where conditions favour high speeds.
Our Power Model analysis explicitly quantifies this: the observed ratio falls
below the model prediction precisely because of this selection effect.
However, selection does not neutralize the risk---it only partially offsets it.
The \SI{42}{\percent} residual gap persists after controlling for the structural
advantages.
Moreover, as Elvik~\cite{Elvik2010} notes, speed choice on these sections
imposes externalities on other road users (trucks, construction vehicles, motorists
with lower maximum speeds) who cannot avoid the risk created by high-speed
drivers.

\paragraph{Comparison with the Netherlands.}
German motorway fatalities per million inhabitants are \SI{40}{\percent}--\SI{50}{\percent}
higher than Dutch motorway fatalities, even though both countries have comparable
vehicle fleets, trauma care systems, and road infrastructure quality.
The Netherlands has operated a maximum motorway speed of \SI{130}{\kilo\metre\per\hour}
(reduced to \SI{100}{\kilo\metre\per\hour} daytime from 2020) under the
Sustainable Safety programme, which explicitly treats speed homogeneity as
a core safety principle~\cite{Wegman2008}.
While confounders exist---notably differences in motorway density and traffic
composition---the magnitude of the gap is difficult to explain without
reference to speed policy.

\paragraph{Limitations.}
This study has three main limitations.
(i)~No vehicle-kilometres-travelled (VKT) normalisation is possible at the
section level from Unfallatlas; the per-accident metrics used here cannot
rule out differences in accident rates (as opposed to severity) driven by
traffic volume.
(ii)~OSM \texttt{maxspeed} tag completeness is user-contributed and may
misclassify some sections.
(iii)~The study period (2016--2024) is too short to detect the trend divergence
with statistical confidence.

%----------------------------------------------------------------------
\section{Conclusion}
%----------------------------------------------------------------------

Using three independent methods applied to the largest available German accident
microdata, we find consistent and mutually reinforcing evidence that
unlimited-speed Autobahn sections are associated with substantially elevated
fatality risk relative to speed-limited sections of the same network.
The key findings are: a \SI{42}{\percent} higher case fatality rate; a 2--2.5$\times$
increase in lethality for the same high-energy collision types; and a Power Model
prediction suggesting the observed gap is a lower bound on the true causal effect.

We conclude that the German \textit{Richtgeschwindigkeit} of \SI{130}{\kilo\metre\per\hour}
(advisory, not mandatory) is insufficient from a Safe System perspective.
A mandatory \SI{130}{\kilo\metre\per\hour} limit, consistent with the rest of the
European Union, would be expected to prevent an estimated 20--40 additional
fatalities per year on currently unlimited sections, in addition to well-documented
reductions in CO\textsubscript{2} emissions and noise.
Given the convergence of evidence, the burden of proof has shifted: it is no longer
tenable to argue that the absence of a speed limit is neutral with respect to
safety.

%----------------------------------------------------------------------
\begin{acknowledgments}
The author thanks the German Federal Statistical Office (Destatis) for open access
to the Unfallatlas microdata, the Dutch Statistics Office (CBS) for motorway
fatality records, and the broader road safety research community whose work is
cited throughout.
\end{acknowledgments}

%----------------------------------------------------------------------
\bibliographystyle{apsrev4-2}
\bibliography{references}

\end{document}
